\section{The positive-entropy law}
\label{sec:positive}

Assume throughout this section that $\lambda=\rho(A)>1$ and
$r=\lambda^{-1}<1$.  Choose positive right and left Perron eigenvectors
\[
 Ax=\lambda x,\qquad y^TA=\lambda y^T,
\]
with no normalization imposed, and define the vertex weight
\begin{equation}
 \pi_v=\frac{x_vy_v}{y^Tx}.
 \label{eq:pi}
\end{equation}
This quantity depends on the chosen presentation and vertex.
All $O$-constants in this section may depend on the fixed pair $(A,v)$.

\begin{lemma}[Derivative normalization]
\label{lem:derivative}
At the old entropy radius,
\begin{equation}
 F_v'(r)=\frac{\lambda}{\pi_v}
 =-\frac{D'(r)}{C_v(r)}.
 \label{eq:derivative}
\end{equation}
\end{lemma}

\begin{proof}
The spectral projector at the simple Perron eigenvalue is
$P=xy^T/(y^Tx)$.  The diagonal resolvent therefore has the local expansion
\begin{equation}
 [(I-zA)^{-1}]_{vv}
 =\frac{\pi_v}{1-\lambda z}+H_v(z),
 \label{eq:resolvent}
\end{equation}
where $H_v$ is analytic near $r$.  Taking the reciprocal in
\cref{eq:resolvent} and using \cref{eq:renewal} gives
\[
 1-F_v(z)=\frac{1-\lambda z}{\pi_v}
 +O\bigl((1-\lambda z)^2\bigr).
\]
Differentiation at $r$ yields the first equality in
\cref{eq:derivative}.  Differentiating
\cref{eq:F-determinant} at $r$, where $D(r)=0$, yields the second.
\end{proof}

\begin{theorem}[Normalized positive-entropy excess]
\label{thm:positive}
For the local split extensions of \cref{thm:construction}, let
\[
 \Gap_N=h(X_{G_N})-h(X_G).
\]
As $N\to\infty$,
\begin{equation}
 \Gap_N=\pi_v r^N\{1+O(Nr^N)\}.
 \label{eq:positive-law}
\end{equation}
Equivalently,
\begin{equation}
 \lim_{N\to\infty}
 \frac{\Gap_N}{\pi_v\lambda^{-N}}=1.
 \label{eq:ratio-limit}
\end{equation}
In particular, for every $\varepsilon>0$, all sufficiently large $N$ give a
nonconjugate split mixing extension with
$0<\Gap_N<\varepsilon$.
\end{theorem}

\begin{proof}
We first show $r_N\to r$.  Fix $s<r$.  Since $F_v(s)<1$ and $s^N\to0$, for
all sufficiently large $N$ one has $F_v(s)+s^N<1$.  The unique root in
\cref{prop:radius} then satisfies $s<r_N<r$.  Letting $s\uparrow r$ proves
the convergence.

Put $\delta_N=r-r_N$.  From
$F_v(r)=1$ and \cref{eq:radius}, the mean-value theorem gives
\begin{equation}
 \delta_NF_v'(\xi_N)=r_N^N
 \label{eq:mvt}
\end{equation}
for some $\xi_N\in(r_N,r)$.  By \cref{lem:derivative} and $r_N\to r$,
$F_v'(\xi_N)$ is bounded away from zero.  Hence
$\delta_N=O(r^N)$ and
\[
 N\delta_N=O(Nr^N)=o(1).
\]
It follows that
\begin{align}
 r_N^N
 &=r^N\left(1-\frac{\delta_N}{r}\right)^N
 =r^N\{1+O(Nr^N)\},
 \label{eq:power-expand}\\
 \frac1{F_v'(\xi_N)}
 &=\frac1{F_v'(r)}\{1+O(r^N)\}.
 \label{eq:derivative-expand}
\end{align}
Combining \cref{eq:mvt,eq:power-expand,eq:derivative-expand},
\[
 r-r_N=\frac{r^N}{F_v'(r)}\{1+O(Nr^N)\}.
\]
Finally,
\begin{align*}
 \Gap_N
 &=\log\frac r{r_N}
 =-\log\left(1-\frac{\delta_N}{r}\right)\\
 &=\frac{r^{N-1}}{F_v'(r)}\{1+O(Nr^N)\}\\
 &=\pi_vr^N\{1+O(Nr^N)\},
\end{align*}
where the last equality is \cref{eq:derivative} and $r\lambda=1$.
This proves \cref{eq:positive-law,eq:ratio-limit}.  Strict positivity also
follows from $\rho(A_N)>\rho(A)$, and convergence to zero gives the final
claim.
\end{proof}

\begin{example}[One-vertex exponent check]
\label{ex:one-vertex}
If $A=[q]$ with $q\ge2$, then $F(z)=qz$ and
\[
 qr_N+r_N^N=1
 \quad\Longleftrightarrow\quad
 \lambda_N=q+\lambda_N^{-(N-1)}.
\]
Here $\pi_v=1$, so \cref{eq:positive-law} reads
\[
 \log\lambda_N-\log q\sim q^{-N}.
\]
The Perron-root equation contains exponent $N-1$, but division by the old
radius in the logarithmic entropy gap restores the exponent $N$.  This
simple case detects a common normalization error.
\end{example}

\begin{remark}
For a nonsymmetric matrix, $\pi_v$ cannot be recovered from an arbitrarily
normalized right eigenvector alone.  The invariant expression
$x_vy_v/(y^Tx)$ is exactly the diagonal entry of the Perron spectral
projector.
\end{remark}
